% experiments.tex — Section 4 for ask-factors NeurIPS 2026 paper
% Depends on: tables, figures in ../figures/

\section{Experiments}\label{sec:experiments}

We evaluate the AFM on a broad universe of US mutual funds spanning 10 asset classes. Our experiments address four questions: (1) Does the portfolio bottleneck autoencoder outperform traditional factor models? (2) How does architecture complexity affect generalization? (3) Does multi-asset training improve performance via cross-asset transfer? (4) Where do the largest gains arise, and where do current limitations remain?

\subsection{Experimental Setup}\label{sec:setup}

\paragraph{Data.}
We construct a dataset of approximately 1,000 US mutual funds drawn from 10 Morningstar categories: large cap, mid cap, small cap, emerging market (EM) equity, non-US equity, US REIT, US aggregate bonds, US high yield debt, EM debt, and municipal debt. Daily excess returns are computed over the risk-free rate. Fund-level characteristics---economic sector holdings (e.g., technology, healthcare, energy) and asset class allocations (e.g., equity, fixed income, cash)---are extracted from regulatory filings (SEC N-PORT) and cross-sectionally normalized per date. This characteristic set constitutes the alternative data that conditions both portfolio weights and factor loadings.

\paragraph{Evaluation metric.}
We use cross-sectional out-of-sample (OOS) $R^2$ as the primary metric. Returns are split temporally: the last 10\% of dates form the held-out evaluation period, ensuring strict separation of training and test data. For each fund $i$, we compute the OOS $R^2$ of the model's fitted returns $\hat{r}_{i,t}$ against realized returns $r_{i,t}$ over the evaluation period. We report the average and median OOS $R^2$ across funds, as well as per-category breakdowns.

\paragraph{Baseline.}
We compare against the Fama-French 4-factor model~\citep{Carhart1997persistence, FamaFrench2015fivefactor}, which includes the market (MKT-RF), size (SMB), value (HML), and momentum (UMD) factors. Factor loadings are estimated via OLS on the training period and applied to the held-out dates. This is the standard benchmark for cross-sectional return explanation and provides a strong baseline on equity categories.

\paragraph{Model configurations.}
We evaluate four AFM variants that vary along two dimensions---architecture and training universe:
\begin{itemize}[leftmargin=*,itemsep=2pt]
    \item \textbf{AFM MLP All}: single-layer MLP backbone (hidden dim 128), trained on all 10 asset classes (34.8M parameters).
    \item \textbf{AFM ResMLP All}: residual MLP backbone (128$\to$64), trained on all 10 asset classes (69.8M parameters).
    \item \textbf{AFM MLP Large Cap}: single-layer MLP, trained on large cap funds only (13.8M parameters).
    \item \textbf{AFM ResMLP Large Cap}: residual MLP, trained on large cap funds only (27.7M parameters).
\end{itemize}
All AFM models extract $F=5$ factors (2 long-only, 3 long-short market-neutral), use a lookback window of $L=252$ trading days, state dimension $S=16$, GELU activations, layer normalization, dropout of 0.3--0.5, and are trained for 150 epochs with Adam (learning rate 0.01, batch size 128). The orthogonality penalty weight ranges from $\lambda_{\text{orth}} = 10^{-5}$ to $10^{-2}$. Full hyperparameter details are provided in Appendix~A.

\paragraph{Compute.}
All experiments are conducted on a single NVIDIA GPU. Each AFM training run completes within several hours. The Fama-French baseline requires only OLS regression and runs in seconds.
% TODO: confirm exact GPU model and wall-clock times

\subsection{Main Results}\label{sec:main_results}

Table~\ref{tab:main_results} presents OOS $R^2$ across all five models and 10 asset classes. The best AFM configuration (MLP backbone, all-asset training) achieves an average OOS $R^2$ of 85.0\%, compared to 68.4\% for the Fama-French 4-factor model---a 16.6 percentage point improvement.

\begin{table}[t]
\centering
\caption{\textbf{Out-of-sample $R^2$ (\%) across models and asset classes.} Each cell reports the average OOS $R^2$ for funds in that category. Higher is better. Bold indicates the best result per category. The AFM MLP All model achieves the highest average $R^2$, with the largest gains on non-equity assets where Fama-French factors have near-zero explanatory power.}
\label{tab:main_results}
\small
\begin{tabular}{lccccc}
\toprule
\textbf{Category} & \textbf{AFM MLP} & \textbf{AFM ResMLP} & \textbf{AFM MLP} & \textbf{AFM ResMLP} & \textbf{FF} \\
 & \textbf{All} & \textbf{All} & \textbf{Large Cap} & \textbf{Large Cap} & \textbf{4-Factor} \\
\midrule
Large Cap       & 90 & 89 & \textbf{95} & 94 & 87 \\
Mid Cap         & \textbf{89} & 86 & 88 & 86 & 85 \\
Small Cap       & 86 & 83 & 84 & 82 & \textbf{88} \\
EM Equity       & 80 & \textbf{83} & 48 & 48 & 43 \\
Non-US Equity   & \textbf{85} & 77 & 69 & 67 & 62 \\
US REIT         & \textbf{66} & 65 & 63 & \textbf{66} & 55 \\
US High Yield   & 53 & \textbf{63} & 47 & 43 & 35 \\
US Agg Bonds    & \textbf{84} & 80 & 11 & 17 & 11 \\
EM Debt         & \textbf{44} & \textbf{44} & 14 & 13 & 11 \\
Municipal Debt  & 27 & \textbf{30} & 5 & 7 & 4 \\
\midrule
\textbf{Average}  & \textbf{85.0} & 82.7 & 73.0 & 70.7 & 68.4 \\
\textbf{Median}   & \textbf{88.1} & 85.4 & 86.0 & 81.5 & 75.4 \\
\bottomrule
\end{tabular}
\end{table}

\paragraph{Equity categories.} All models perform well on equity categories. Large cap, mid cap, and small cap $R^2$ exceeds 82\% for every model, including Fama-French. This is expected: the Fama-French factors are designed for US equities and capture the dominant sources of variation in these categories. Notably, the Fama-French model achieves 88\% on small cap, slightly outperforming the AFM MLP All (86\%), which may reflect the tight alignment between the SMB factor and small-cap return variation.

\paragraph{Non-equity categories.} The AFM's advantages emerge most clearly on asset classes where traditional equity factors have limited explanatory power. On US aggregate bonds, the AFM MLP All achieves 84\% $R^2$ compared to just 11\% for Fama-French---a 73 percentage point gap. Similarly, on EM debt (44\% vs.\ 11\%), US high yield (53\% vs.\ 35\%), and US REITs (66\% vs.\ 55\%), the AFM substantially outperforms. These categories are driven by interest rate, credit, and currency risk factors that equity-based models cannot capture. The portfolio bottleneck learns cross-asset factor portfolios that span these risk dimensions.

Figure~\ref{fig:r2_comparison} provides a visual comparison across all categories and models, highlighting the systematic advantage of multi-asset AFM training on non-equity categories.

\begin{figure}[t]
\centering
\includegraphics[width=\linewidth]{figures/r2_comparison}
\caption{\textbf{OOS $R^2$ by asset class and model.} Grouped bar chart comparing all five models across 10 mutual fund categories. The AFM models trained on all asset classes (blue and orange bars) dominate on non-equity categories (right side), while all models perform comparably on equity categories (left side). Fama-French (gray) achieves near-zero $R^2$ on bonds and EM debt, where the AFM captures interest rate and credit factors through its portfolio bottleneck.}
\label{fig:r2_comparison}
\end{figure}

\subsection{Ablation Studies}\label{sec:ablation}

We conduct two ablations to isolate the contributions of architecture choice and training universe composition. Table~\ref{tab:ablation} summarizes the results in a $2 \times 2$ grid.

\begin{table}[t]
\centering
\caption{\textbf{Ablation: architecture and training universe.} Each cell reports average OOS $R^2$ (\%) and parameter count. The single-layer MLP consistently outperforms the deeper residual MLP, and multi-asset training (``All'') consistently outperforms large-cap-only training. Higher is better.}
\label{tab:ablation}
\begin{tabular}{lcc}
\toprule
 & \textbf{All Assets} & \textbf{Large Cap Only} \\
\midrule
\textbf{MLP} (single layer)     & \textbf{85.0\%} (34.8M) & 73.0\% (13.8M) \\
\textbf{Residual MLP} (2 layers) & 82.7\% (69.8M) & 70.7\% (27.7M) \\
\bottomrule
\end{tabular}
\end{table}

\paragraph{Architecture complexity.}
The single-layer MLP (34.8M parameters) outperforms the deeper residual MLP (69.8M parameters) by 2.3 percentage points in the all-asset setting (85.0\% vs.\ 82.7\%) and by 2.3 percentage points in the large-cap setting (73.0\% vs.\ 70.7\%). This result is consistent across both training universes, suggesting that the additional depth and skip connections of the residual architecture add capacity without useful expressiveness for this task. We hypothesize that mutual fund returns exhibit smoother cross-sectional structure than individual stock returns---funds are diversified portfolios, so their return cross-section is lower-dimensional---and a single hidden layer suffices to capture the relevant nonlinearities in the characteristics-to-loadings mapping.

\paragraph{Training universe.}
Multi-asset training yields a 12.0 percentage point improvement over large-cap-only training for the MLP backbone (85.0\% vs.\ 73.0\%) and a 12.0 percentage point improvement for the residual MLP (82.7\% vs.\ 70.7\%). The gains are concentrated in non-equity categories: US aggregate bonds improve from 11\% to 84\%, EM debt from 14\% to 44\%, and non-US equity from 69\% to 85\%. Even equity categories benefit modestly---EM equity improves from 48\% to 80\%. This demonstrates cross-asset transfer learning: by training jointly across equities, bonds, REITs, and emerging markets, the model discovers shared factor structure (e.g., common interest rate and credit risk drivers) that single-asset training cannot identify. The portfolio bottleneck is essential to this transfer---it forces the model to construct factor portfolios whose returns are informative for \emph{all} asset classes simultaneously.

\subsection{Cross-Asset Analysis}\label{sec:cross_asset}

To better understand where the AFM's gains originate, we group the 10 categories by asset type. \emph{Equity categories} (large cap, mid cap, small cap, EM equity, non-US equity) achieve above 80\% $R^2$ for all AFM all-asset models, with the AFM and Fama-French performing comparably on domestic equities. The differentiation between models comes primarily from three groups:

\begin{itemize}[leftmargin=*,itemsep=2pt]
    \item \textbf{International and EM equity:} The AFM captures diversification and currency factors that US-centric models miss. EM equity improves from 43\% (FF) to 80--83\% (AFM All).
    \item \textbf{Fixed income:} US aggregate bonds (84\% vs.\ 11\%) and US high yield (53--63\% vs.\ 35\%) show the largest absolute gains. The portfolio bottleneck learns rate-sensitive and credit-sensitive factor portfolios that have no analogue in the Fama-French framework.
    \item \textbf{Alternative and EM debt:} US REITs (66\% vs.\ 55\%) and EM debt (44\% vs.\ 11\%) benefit from cross-asset training but remain below equity-level $R^2$, reflecting the distinct risk drivers and thinner cross-sections in these categories.
\end{itemize}

Municipal debt remains the most challenging category, with $R^2$ peaking at 30\%. We discuss this limitation in Section~\ref{sec:limitations}.

\subsection{Training Dynamics}\label{sec:training_dynamics}

Figure~\ref{fig:training_loss} shows the training and validation loss trajectories for all four AFM variants. All models converge rapidly within 10--15 epochs, with continued modest improvement thereafter. The best validation loss is reached between epochs 84 and 114, depending on the configuration. Importantly, no overfitting is observed for any model: validation loss stabilizes and does not deteriorate even after 150 epochs of training. This suggests that the portfolio bottleneck acts as an effective implicit regularizer---by constraining the latent factors to be returns of explicit portfolios rather than free latent codes, the model's representational capacity is disciplined in a way that prevents memorization of training-period return patterns.

\begin{figure}[t]
\centering
\includegraphics[width=\linewidth]{figures/loss_trajectories}
\caption{\textbf{Training loss trajectories for all AFM variants.} Epoch (x-axis) vs.\ reconstruction loss (y-axis) for training (solid) and validation (dashed) sets. All four configurations---MLP All (blue), ResMLP All (orange), MLP Large Cap (green), ResMLP Large Cap (red)---converge within 10--15 epochs and show no signs of overfitting through 150 epochs. Best validation loss is reached between epochs 84--114.}
\label{fig:training_loss}
\end{figure}

\subsection{Limitations}\label{sec:limitations}

Several limitations of the current evaluation deserve emphasis. First, \textbf{municipal debt} $R^2$ peaks at 30\%, likely because our fund-level characteristics (sector holdings and asset class allocations) do not capture the credit quality, duration, and tax-status features that drive municipal bond returns. Richer characteristic inputs---including credit ratings, effective duration, and state-level tax information---would likely improve performance on this category. Second, our evaluation is based on \textbf{explanatory $R^2$} only; we do not assess portfolio-level performance metrics such as Sharpe ratios, turnover, or transaction costs, which are critical for practical deployment. Third, we report results on a \textbf{single temporal split} (last 10\% of dates) rather than rolling or cross-validated evaluation, due to the computational cost of training. While the temporal split is standard practice in asset pricing~\citep{GuKellyXiu2020empirical, GuKellyXiu2021autoencoder}, future work should examine robustness across different evaluation windows.
